(1) 牛顿方程 \(m\ddot u_n=K_1(v_n-u_n)+K_2(v_{n-1}-u_n)\)，\(m\ddot v_n=K_1(u_n-v_n)+K_2(u_{n+1}-v_n)\)。代入 Bloch 波得矩阵 \(\begin{pmatrix}K_1+K_2-m\omega^2&-(K_1+K_2e^{-ika})\\-(K_1+K_2e^{ika})&K_1+K_2-m\omega^2\end{pmatrix}\binom UV=0\)。非零振幅需行列式为零，复相位来自胞间弹簧且保持 Hermitian。

(2) 行列式方程 \((K_1+K_2-m\omega^2)^2=|K_1+K_2e^{ika}|^2=K_1^2+K_2^2+2K_1K_2\cos ka\)。故 \(m\omega_\pm^2=K_1+K_2\pm\sqrt{K_1^2+K_2^2+2K_1K_2\cos ka}\)，减号为声学支，加号为光学支。两支均非负，这还可由弹簧势能平方和直接验证。

(3) 当 \(k=0\)，根号为 \(K_1+K_2\)，得到 \(\omega_-(0)=0,\omega_+^2(0)=2(K_1+K_2)/m\)。在 \(k=\pi/a\)，根号为 \(|K_1-K_2|\)，故 \(\omega_-^2=2\min(K_1,K_2)/m,\omega_+^2=2\max(K_1,K_2)/m\)。若 \(K_1\ne K_2\)，两个频段中间存在无传播 Bloch 正则模的频隙；两弹簧相同时禁带闭合。

(4) 令 \(ka\ll1\)，\(\cos ka=1-(ka)^2/2+O(k^4)\)，根号展开为 \((K_1+K_2)-K_1K_2(ka)^2/[2(K_1+K_2)]+\cdots\)，于是 \(\omega_-(k)\simeq a\sqrt{K_1K_2/[2m(K_1+K_2)]}|k|\)，声速 \(c=a\sqrt{K_1K_2/[2m(K_1+K_2)]}\)。当 \(K_1=K_2=K\)，实际平移周期缩为 \(a/2\)，双支表示只是同一单原子链色散在较小布里渊区折叠而成。
